\clearpage
\VerdierPhysicalPageAnchor{0068}%
\VerdierSemanticAnchor{ch1-2-2-1-continuation}%
$j<i$ ならば $\varepsilon(e_j,i)=-1$、$i\leq j$ ならば
$\varepsilon(e_j,i)=1$ であることが直ちに分かる。さらに、写像
$(\sigma,i)\mapsto\varepsilon(\sigma,i)$ はこれらの性質によって一意に
定まる。ここから関係式

\VerdierSemanticAnchor{ch1-2-2-1-eq2}
\begin{equation}
  \varepsilon(\sigma,i)\varepsilon(\sigma+e_i,j)
  =-\varepsilon(\sigma,j)\varepsilon(\sigma+e_j,i)
  \tag{2.2.1.2}
\end{equation}
が直ちに従う。これは各 $\sigma\in\mathbb Z[n]$ および
$(i,j)\in[n]\times[n]$、$i\neq j$ に対して成り立つ。

\VerdierSemanticAnchor{ch1-2-2-2}
\noindent\textbf{2.2.2.}\quad
$Y$ を $C$ における $n$-重複複体とする。次のようにおく：

\VerdierSemanticAnchor{ch1-2-2-2-eq1}
\begin{equation}
  \delta_Y^{\sigma,i}=\varepsilon(\sigma,i)d_Y^{\sigma,i}.
  \tag{2.2.2.1}
\end{equation}

\VerdierSemanticAnchor{ch1-2-2-2-eq2}
(2.2.1.2) から、関係式 (2.1.2.1) は次の関係式と同値である：
\begin{equation}
  \begin{aligned}
    \delta_Y^{\sigma+e_i,i}\delta_Y^{\sigma,i}
      &=0,
      &&\forall\sigma\in\mathbb Z[n],\quad\forall i\in[n],\\
    \delta_Y^{\sigma+e_j,i}\delta_Y^{\sigma,j}
      +\delta_Y^{\sigma+e_i,j}\delta_Y^{\sigma,i}
      &=0,
      &&\forall\sigma\in\mathbb Z[n],\quad\forall i,j\in[n].
  \end{aligned}
  \tag{2.2.2.2}
\end{equation}

\VerdierSemanticAnchor{ch1-2-2-2-pushforward}
ここで $\varphi\colon[n]\to[m]$ を写像とし、$Y$ が
$\varphi$-総和可能であるとする。次の公式によって $m$-重複複体
$\int_\varphi Y$ を定める：

\VerdierSemanticAnchor{ch1-2-2-2-eq3}
\begin{equation}
  \left(\int_\varphi Y\right)^\tau
  =\bigoplus_{\varphi(\sigma)=\tau}Y^\sigma,
  \qquad \forall\tau\in\mathbb Z[m],
  \tag{2.2.2.3}
\end{equation}

\VerdierSemanticAnchor{ch1-2-2-2-eq4}
\begin{equation}
  \begin{aligned}
    \delta_{\int_\varphi Y}^{\tau,i}(y_\sigma)
      &=\sum_{\varphi(j)=i}\delta_Y^{\sigma,j}(y_\sigma),\\[-1mm]
    &\hspace{-18mm}
      \sigma\in\mathbb Z[n],\quad i\in[m],\quad
      y_\sigma\in Y^\sigma,\quad\tau=\varphi(\sigma).
  \end{aligned}
  \tag{2.2.2.4}
\end{equation}

\VerdierSemanticAnchor{ch1-2-2-2-additive-interpretation}
$C$ がアーベル群の圏であるとき、公式 (2.2.2.4) は射
\[
  \left(\int_\varphi Y\right)^\tau
  =\bigoplus_{\varphi(\sigma)=\tau}Y^\sigma
  \longrightarrow
  \left(\int_\varphi Y\right)^{\tau+e_i}
  =\bigoplus_{\varphi(\sigma)=\tau+e_i}Y^\sigma
\]
を定める。
$C$ が加法圏であるときも、$C$ 上のアーベル群値反変関手の圏で考えること、
すなわち $y_\sigma$ を $C$ の任意の対象から $Y^\sigma$ への射と解釈する
ことによって、$\left(\int_\varphi Y\right)^\tau$ から
$\left(\int_\varphi Y\right)^{\tau+e_i}$ への射を定められる。

\VerdierSemanticAnchor{ch1-2-2-2-verification}
直ちに計算すれば、$\delta_{\int_\varphi Y}^{\tau,i}$ は関係式
(2.2.2.2) を満たす。したがって微分に戻すと、(2.2.2.3) で定めた対象に
次の射を備えたものは

\VerdierSemanticAnchor{ch1-2-2-2-eq5}
\begin{equation}
  \begin{aligned}
    d_{\int_\varphi Y}^{\tau,i}(y_\sigma)
      &=\varepsilon(\tau,i)
        \sum_{\varphi(j)=i}
          \varepsilon(\sigma,j)d_Y^{\sigma,j}(y_\sigma),\\[-1mm]
    &\hspace{-18mm}
      \tau\in\mathbb Z[m],\quad i\in[m],\quad
      \sigma\in\mathbb Z[n],\quad\varphi(\sigma)=\tau
  \end{aligned}
  \tag{2.2.2.5}
\end{equation}
